\section{Model, prior work, and disaggregation}\label{sec:foundations}

\subsection{The sparse network--simplex hull}\label{subsec:model}
Let $D=(V,E)$ be a finite directed multigraph. Parallel arcs and self-loops
are permitted. Its incidence matrix $A\in\{-1,0,1\}^{V\times E}$ has a $-1$
at the tail and a $+1$ at the head of a non-loop arc; a self-loop has a zero
column. Thus $Ax=b$ uses the convention \emph{incoming minus outgoing}.
For $b\in\Q^V$ and finite $u\in\Q_+^E$, set
\begin{equation}\label{eq:flow-polytope}
 \Flow=\{x\in\R^E:Ax=b,\ 0\le x\le u\}.
\end{equation}
Lower bounds can be handled by translating $x$ and the corresponding product
coordinates. Throughout, graph assumptions concern the graph in
\eqref{eq:flow-polytope} after any such modeling transformations.

For an integer $m\ge0$, write $[m]=\{1,\ldots,m\}$ and
\begin{equation}\label{eq:simplex}
 \Delta_m=\{y\in\R_+^m:\one^\top y\le1\},\qquad
 \lambda_j=y_j\ (j\in[m]),\qquad
 \lambda_0=1-\one^\top y.
\end{equation}
The \emph{states} are the $d=m+1$ vertices $e^0=0,e^1,\ldots,e^m$ of
$\Delta_m$; $e^j$ is the $j$th unit vector when $j\ge1$. The residual state
is state $0$. An observation set $\Obs\subseteq E\times[m]$ specifies
the products retained as original coordinates. We study
\begin{align}
 S&=\{(x,y,z):x\in\Flow,\ y\in\Delta_m,
             \ z_{ej}=x_e y_j\quad((e,j)\in\Obs)\},\label{eq:graph-set}\\
 \Hull&=\conv S.\label{eq:sparse-hull}
\end{align}
An unobserved product is absent from $z$, even if it is introduced temporarily
in an extended formulation. State $0$ has no observed coordinates.
We use $\conv\varnothing=\varnothing$. When $m=0$, $y$ and $z$ are empty
vectors and $\Hull=\Flow$.

The original coordinates are $(x,y,z)$. An extended formulation is a linear
system in these coordinates and additional continuous variables whose
projection is $\Hull$. A separation procedure must either certify membership
or produce a linear inequality valid for all of $\Hull$ and strictly violated
by the candidate. Unless specified otherwise, a facet is taken relative to
the affine hull. A bound on integer coefficients presupposes a stated
normalization; a ratio of two nonzero coefficients is invariant under scaling
the inequality. Arithmetic-operation bounds are distinguished from bounds
on rational encoding length.

For comparison, the usual independent product relaxation adds, for every
$(e,j)\in\Obs$, the four inequalities
\begin{equation}\label{eq:mccormick}
 0\le z_{ej},\qquad z_{ej}\le u_e y_j,\qquad
 z_{ej}\le x_e,\qquad
 z_{ej}\ge x_e-u_e(1-y_j)
\end{equation}
to $x\in\Flow$ and $y\in\Delta_m$. These are the McCormick inequalities
for the box $[0,u_e]\times[0,1]$ \citep{McCormick1976}. They need not describe
the joint hull in \eqref{eq:sparse-hull}.

\subsection{Relation to prior work}\label{subsec:prior-work}
Disjunctive programming \citep{Balas1998} over the simplex vertices and
the simultaneous convexification of \citet{DavarniaRichardTawarmalani2017},
which covers linear combinations of bilinear products over a polytope times
a simplex, already give the hull. Sparsity of $\Obs$ therefore does not make
the existence of an exact extended formulation new.

The closest predecessor is \citet{KhademniaDavarnia2025}. Their Theorem~1
gives a complete original-coordinate hull description through all extended
cancel-and-relax (EC\&R) inequalities. Their Theorem~2 gives explicit tree
constructions for one simplex variable, and their Theorem~3 a structured
forest construction for multiple simplex variables; that forest recipe has
narrower scope than their complete projection framework. Their Section~3
represents equality balances by opposite inequalities, so their setting
includes \eqref{eq:flow-polytope}. We use the equality circulation space to
refine this established hull: residual-cycle counts, complete finite
libraries for the specified graph classes, and sharp coefficient boundaries.

Two established reductions are especially close to our formulation.
\citet[Proposition~2.6]{Davarnia2016} proves that convexifying Cartesian
components separately is exact when they share a simplex. We apply that
gluing principle to the independent block circulation spaces and explicitly
refine the blockwise merged states into common global states, including
zero weights. \citet[Theorem~3.1]{LibertiPantelides2006} use multiplied linear
equations to eliminate common-factor bilinear terms. Our nullity count is
the graph case of that linear algebra: the undetermined state flows are
exactly the circulations on unobserved arcs. Our formulation also keeps all
state and residual capacity bounds and controls the coefficients after
elimination. This matters because eliminating a bilinear equation does not
by itself justify discarding its convex relaxation, as
\citet[discussion after Theorem~3.1]{LibertiPantelides2006} observe.

For direct descriptions, \citet{KisHorvath2022} represent Cayley hulls by
networks and derive projection inequalities from cuts. Their Sections~5.7
and~5.9 are close predecessors of our parallel-path construction; in
particular, Proposition~22 and equations~(30)--(31) in Section~5.9 use a
lower-bound shift and a transportation projection. We specialize these
mechanisms to observed state slices and give the resulting separator and
recovery. For bounded-rank blocks, the support geometry of Minkowski sums
and refinements by edge-generated zonotopes is classical
\citep{GritzmannSturmfels1993,OnnRothblum2004}. Our bounds use the
totally unimodular cycle matrix to control all state directions, support
normals, and recovery bases, including degenerate slices.

Coefficient growth also has precedents. In the published article,
\citet[Example~2]{KhademniaDavarnia2025} already exhibit a nonunit dual
aggregation weight with two simplex variables. Such a multiplier alone need
not force a nonunit ratio between product coefficients in every
original-coordinate description. Our coordinate-section arguments prove
this stronger conclusion for the stated restricted networks, and it is
invariant under adding affine-hull equations. Large coefficients for $0/1$ polytopes
\citep{AlonVu1997} and transfers of facets to graph polytopes
\citep{Fiorini2006} are established; our restrictions concern the actual
retained product coordinates, the label count, and the graph structure.

Likewise, \citet[Theorems~1.1 and~1.2]{DeLoeraOnn2006} establish
three-way transportation and bitransportation universality, including a
two-commodity network interpretation. Our unrestricted-network result
transfers that universality to a coordinate section of a sparse bilinear
hull with two explicit simplex labels and unit network data. On
series--parallel networks, integrality of aggregate multiflows does not
settle the question: \citet[Remark~1]{AlmoghrabiSkutellaWarode2026}
distinguish aggregate integrality from decomposability of the individual
commodity-flow vector. Observed products retain individual-state
information that aggregate flows omit.

\citet{DavarniaRahimian2026} develop aggregation and projection for
bilinear constraints with binary simplex variables, with applications to
mixing sets in chance-constrained programming. Their target is a hull in
the $x$-space after eliminating the simplex variables (Section~3). Here $y$
is continuous and remains in the hull together with the selected products.

\begin{table}[tb]\centering\small
\caption{Established ingredients and the specific developments in this paper.
The comparison concerns the stated constructions and bounds, not priority
for general convexification or polynomial-time separation.}
\label{tab:prior-comparison}
\begin{tabular}{>{\raggedright\arraybackslash}p{.43\linewidth}
                >{\raggedright\arraybackslash}p{.49\linewidth}}\toprule
Established ingredient & Development here\\\midrule
Simplex disaggregation and complete EC\&R projection
\citep{DavarniaRichardTawarmalani2017,KhademniaDavarnia2025}
& Exact residual-cycle formulation with unit coefficients;
\cref{thm:observed-compression}.\\[4pt]
Shared-simplex gluing and elimination by multiplied equalities
\citep{Davarnia2016,LibertiPantelides2006}
& Blockwise observation merging and recovery, with
$\sum_{B,j}\rho_{Bj}$ residual coordinates and the individual-product
completion interpretation; \cref{thm:block-state-reduction,prop:minimum-completion}.\\[4pt]
Transportation projection and support refinements
\citep{KisHorvath2022,GritzmannSturmfels1993,OnnRothblum2004}
& Complete graph-specific libraries and coefficient bounds;
\cref{thm:bounded-rank,thm:rank-recovery,cor:chain-observed-labels}.\\[4pt]
Nonunit aggregation weights, large polytope coefficients,
and transportation universality
\citep{KhademniaDavarnia2025,AlonVu1997,DeLoeraOnn2006}
& Unavoidable product-coefficient ratios: the sharp flat-chain threshold,
sparse Fibonacci family, and two-label coordinate-section transfer;
\cref{thm:three-state-chain,cor:fibonacci-simple,thm:universality}.\\\bottomrule
\end{tabular}
\end{table}

\subsection{The known exact disaggregation}\label{subsec:disaggregation}
For $\lambda\ge0$, define the scaled flow system
\begin{equation}\label{eq:scaled-flow}
 \Flow(\lambda)=\{f\in\R^E:Af=\lambda b,\ 0\le f\le\lambda u\}.
\end{equation}
For $\lambda>0$, this is $\lambda\Flow$; at $\lambda=0$ it is $\{0\}$,
including when $\Flow$ is empty. Using \eqref{eq:scaled-flow} avoids assigning
an ambiguous meaning to $0\varnothing$.

\begin{proposition}[Simplex-vertex disaggregation]\label{prop:disaggregation}
A point $(x,y,z)$ belongs to $\Hull$ if and only if $y\in\Delta_m$ and
there exist state flows $f^0,\ldots,f^m\in\R^E$ such that
\begin{equation}\label{eq:disaggregation}
 f^k\in\Flow(\lambda_k)\quad(k\in\states),\qquad
 \sum_{k=0}^m f^k=x,\qquad
 f^j_e=z_{ej}\quad((e,j)\in\Obs).
\end{equation}
The equivalence holds even if $\Flow$ is empty. If the system is feasible,
it provides a convex decomposition into at most $m+1$ points of $S$.
\end{proposition}
\begin{proof}
For a point of $S$, set $f^k=\lambda_k x$. This satisfies all displayed
relations, since $\sum_k\lambda_k=1$. The system is linear in
$(x,y,z,f)$, so its projection is convex and contains $\Hull$.

Conversely, suppose the system is feasible. If $\lambda_k=0$, the finite
upper bounds give $f^k=0$. For each positive weight let
$x^k=f^k/\lambda_k\in\Flow$ and define the graph point
$p^k=(x^k,e^k,z^k)$ by
$z^k_{ej}=x^k_e(e^k)_j$. Its weighted average has $x$ coordinate
$\sum_k f^k=x$, $y$ coordinate $\sum_k\lambda_k e^k=y$, and observed
product coordinate $\lambda_j x^j_e=f^j_e=z_{ej}$ when
$\lambda_j>0$. If $\lambda_j=0$, that observed coordinate is zero on both
sides. Thus $(x,y,z)=\sum_{k:\lambda_k>0}\lambda_k p^k\in\Hull$.
There is at least one positive weight, so feasibility also implies
$\Flow\ne\varnothing$. This proves the empty-set assertion.
\end{proof}

The formulation has $(m+1)|E|$ state-flow variables. Its rational size is
polynomial in the network and simplex input, so polynomial-time optimization
and separation over $\Hull$ already follow from linear programming. The
structural questions are how to avoid this full state-by-arc expansion and
how to project it into useful original-coordinate inequalities.
Eliminating $f^0=x-\sum_{j=1}^m f^j$ leaves the residual bounds
\begin{equation}\label{eq:residual-bounds}
 0\le x-\sum_{j=1}^m f^j\le
       \Bigl(1-\sum_{j=1}^m y_j\Bigr)u.
\end{equation}
These bounds cannot in general be discarded.

\begin{corollary}[Integral arc flows give the same hull]\label{cor:integral-flows}
If $b$ and $u$ are integral, then
\[
 \Hull=\conv\{(x,y,z)\in S:x\in\Z^E\}.
\]
\end{corollary}
\begin{proof}
The bounded flow polytope $\Flow$ is integral, a classical consequence of
network total unimodularity \citep{HoffmanKruskal1956}. Here is a direct
verification covering multigraphs. At a nonintegral feasible flow, consider
the arcs whose values are nonintegral. If there is no loop, each incident
vertex has degree at least two in this subgraph: a vertex with just one
such arc would equate a noninteger plus integers to an integer in its
balance equation. A nonempty subgraph therefore contains a loop or an
undirected cycle. The signed
circulation on that cycle can be added and subtracted in sufficiently small
amounts while preserving the bounds, because each of its arc values is
strictly between its integral bounds. Thus the flow is not a vertex.
All vertices of $\Flow$ are integral, and boundedness expresses every
flow as a convex combination of them.

For a point of $\Hull$, apply \cref{prop:disaggregation}. Decompose each
positive-weight normalized flow into integral vertices of $\Flow$ and
pair those vertices with its fixed simplex vertex. At a fixed simplex
vertex, all selected products are linear in the flow, so these refined
graph points have the required weighted average. This proves the inclusion
in the displayed right-hand side; the reverse inclusion follows from its
containment in $\conv S$. If $\Flow$ is empty both sides are empty.
\end{proof}

The corollary is a hull equality, not a bound on the size of integral
decompositions. Refining a normalized state flow can require several
integral vertices, so the at-most-$m+1$ bound in \cref{prop:disaggregation}
applies to continuous graph points; already for $m=0$, a fractional flow is
not a single integral point. Nor does the corollary assert integrality of
fractional coordinate sections of $\Hull$.

\section{Independent blocks and locally merged states}\label{sec:blocks}

\subsection{Circulation factorization}\label{subsec:block-factorization}
All graph connectivity terminology below refers to the undirected multigraph
underlying $D$. Isolated vertices impose only their balance equation. Let
$E_{\rm br}$ be the bridges. Let $\Blocks$ be the cyclic blocks: the maximal
biconnected edge blocks that contain a cycle, with each self-loop treated as
its own block. A pair of parallel edges can form a cycle of length two.
The sets $E_{\rm br}$ and $E_B$ for $B\in\Blocks$ partition $E$.
The cycle rank of a connected block is
\begin{equation}\label{eq:cycle-rank}
 r_B=|E_B|-|V_B|+1.
\end{equation}
In particular, a loop has rank one. Rank refers to the ambient circulation
space, even when capacities reduce the dimension of the feasible polytope.

The equation $Av=b$ is solvable exactly when the entries of $b$ sum to zero
on every connected component. Necessity follows by summing the incidence
rows in a component. For sufficiency, choose a spanning forest, set the
nonforest coordinates to zero, and determine each tree-edge flow by the
balance sum on one side of its fundamental cut, with the sign determined by
the arc orientation. This satisfies every balance equation. This construction
uses $O(|V|+|E|)$ arithmetic operations, and for rational $b$ produces a
rational $v$ of polynomial encoding length. The reference flow $v$ need not
obey its capacity bounds. If the component condition fails, $\Flow$ and
$\Hull$ are empty.

Fix such a reference vector $v$. For each $B$, let $A_B$ be the incidence
matrix restricted to columns in $E_B$ and rows in $V_B$, and define
\begin{equation}\label{eq:block-domain}
 K_B=\{h\in\R^{E_B}:A_Bh=0,\ -v_{E_B}\le h\le u_{E_B}-v_{E_B}\}.
\end{equation}
Unlike $\Flow$, this is a domain of circulation \emph{deviations}; it need
not contain the origin.

\begin{lemma}[Block factorization]\label{lem:block-factorization}
Every $g\in\ker A$ vanishes on $E_{\rm br}$ and restricts to a circulation
on each cyclic block. Conversely, combining arbitrary block circulations
and zero bridge coordinates gives an element of $\ker A$. Consequently,
if every bridge satisfies $0\le v_e\le u_e$, the map
\begin{equation}\label{eq:block-product-map}
 x_e=v_e\quad(e\in E_{\rm br}),\qquad
 x_{E_B}=v_{E_B}+h_B\quad(B\in\Blocks)
\end{equation}
is an affine bijection from $\prod_{B\in\Blocks}K_B$ to $\Flow$.
If a bridge bound fails, $\Flow$ is empty. The product of an empty family
of block domains is a singleton.
\end{lemma}
\begin{proof}
Choose a spanning forest of the underlying multigraph, omitting loops. Each
nonforest edge gives a signed fundamental-cycle circulation: use the unique
forest path between its ends, or just the edge itself for a loop. These
vectors form a basis of $\ker A$. Indeed, subtracting from a circulation
the combination that cancels all nonforest coordinates leaves a circulation
supported on a forest, which is zero by successive leaf elimination.

Each fundamental cycle lies in a single cyclic block. Hence its vector is
zero on every bridge, and on every block it is either zero or a circulation
of that block. The basis expansion proves the first assertion. For the
converse, extend each block circulation by zero to $E$. Its incidence is
zero at every vertex, including articulation vertices, so their sum is a
global circulation. The sum is direct because the block edge sets are
disjoint. Applying this decomposition to $x-v$ proves that $Ax=b$ is
equivalent to the circulation conditions and the fixed bridge coordinates
in \eqref{eq:block-product-map}. The bounds then separate by edge block,
giving precisely \eqref{eq:block-domain} and the bridge conditions.
\end{proof}

In particular, an articulation vertex needs no coupling variable between
block deviations: its net balance is carried by $v$, and every block
deviation has zero net balance there.
Any spanning tree in $B$ gives a fundamental-cycle matrix
$C_B\in\{0,\pm1\}^{E_B\times r_B}$ with chord rows equal to the identity.
Thus $h_B=C_Bt_B$ gives unique cycle coordinates $t_B\in\R^{r_B}$.
These coordinates remain useful when $K_B$ is lower-dimensional.

\subsection{The state merger and its exact refinement}\label{subsec:state-merger}
For a nonempty convex set $K$, use the usual scalar multiplication
$\alpha K=\{\alpha h:h\in K\}$, so $0K=\{0\}$. For nonnegative scalars,
\begin{equation}\label{eq:homothetic-merger}
 \alpha_1K+\cdots+\alpha_qK
       =(\alpha_1+\cdots+\alpha_q)K.
\end{equation}
If the total weight is positive, the inclusion from left to right follows
by dividing by that weight and taking a convex combination; the reverse
inclusion uses the same point of $K$ in each summand. If the total is zero,
both sides are $\{0\}$. An empty sum is also interpreted as $\{0\}$.

For each block define its observed labels and their count by
\begin{equation}\label{eq:observed-labels}
 J_B=\{j\in[m]: (e,j)\in\Obs\text{ for some }e\in E_B\},
 \qquad a_B=|J_B|.
\end{equation}
This is a structural definition: a label in $J_B$ remains present when
$y_j=0$. Let $*_B$ denote a local merged state, distinct from the global
residual state $0$, and set
\begin{equation}\label{eq:merged-weight}
 I_B=J_B\cup\{*_B\},\qquad
 \lambda_{B,j}=y_j\ (j\in J_B),\qquad
 \lambda_{B,*}=1-\sum_{j\in J_B}y_j.
\end{equation}
The merged state collects all labels absent from $J_B$ and state $0$.
Different blocks may merge different global states.

\begin{theorem}[Exact blockwise state reduction]\label{thm:block-state-reduction}
Assume $\Flow\ne\varnothing$. A point $(x,y,z)$ belongs to $\Hull$ if and
only if $x\in\Flow$, $y\in\Delta_m$, every observed bridge satisfies
\begin{equation}\label{eq:bridge-product}
 z_{ej}=v_e y_j\qquad(e\in E_{\rm br},\ (e,j)\in\Obs),
\end{equation}
and for every cyclic block $B$ there exist vectors
$h^{B,k}\in\R^{E_B}$, $k\in I_B$, satisfying
\begin{align}
 h^{B,k}&\in\lambda_{B,k}K_B\quad(k\in I_B),\label{eq:local-state-domain}\\
 \sum_{k\in I_B}h^{B,k}&=x_{E_B}-v_{E_B},\label{eq:local-state-sum}\\
 h^{B,j}_e&=z_{ej}-y_jv_e
       \quad((e,j)\in\Obs,\ e\in E_B).\label{eq:local-observation}
\end{align}
The block conditions yield state flows in
\eqref{eq:disaggregation}, and hence a decomposition into at most $m+1$
graph points, by explicit proportional refinement.
\end{theorem}
\begin{proof}
By \cref{lem:block-factorization}, every $K_B$ is nonempty. Start with
state flows in \eqref{eq:disaggregation}. Their deviations
$g^k=f^k-\lambda_k v$ are circulations. On $B$, they satisfy
$g^k_{E_B}\in\lambda_kK_B$: division proves this for positive weights,
and both deviations and scaled domains are zero for zero weights.
For $j\in J_B$ take $h^{B,j}=g^j_{E_B}$, and sum the other deviations
into $h^{B,*}$. Identity \eqref{eq:homothetic-merger} proves
\eqref{eq:local-state-domain}. Summing the deviations proves
\eqref{eq:local-state-sum}, and the observed coordinates give
\eqref{eq:local-observation}. On bridges the deviations vanish, proving
\eqref{eq:bridge-product}.

Conversely, fix feasible local vectors. For each global state $k$ and block
$B$, if $k\in J_B$, set $g^k_{E_B}=h^{B,k}$. For the remaining states use
\begin{equation}\label{eq:proportional-refinement}
 g^k_{E_B}=\begin{cases}
 (\lambda_k/\lambda_{B,*})h^{B,*},&\lambda_{B,*}>0,\\
 0,&\lambda_{B,*}=0.
 \end{cases}
\end{equation}
If the merged weight is zero, every constituent weight is zero as well.
In every case $g^k_{E_B}\in\lambda_kK_B$, and the refined vectors sum to
the block deviation. Put $g^k_e=0$ on bridges and combine blocks. The block
factorization gives $Ag^k=0$. Then $f^k=\lambda_kv+g^k$ satisfies
$Af^k=\lambda_kb$. Its capacity bounds hold on blocks by the scaled
domains, and on bridges because $\Flow\ne\varnothing$ implies
$0\le v_e\le u_e$. Equations \eqref{eq:local-state-sum} and
\eqref{eq:local-observation}, together with the bridge equations, give
$\sum_kf^k=x$ and all observed products. Apply
\cref{prop:disaggregation}.
\end{proof}

The domain condition \eqref{eq:local-state-domain} is itself linear:
\begin{equation}\label{eq:linear-scaled-block}
 A_Bh^{B,k}=0,\qquad
 -\lambda_{B,k}v_{E_B}\le h^{B,k}
       \le\lambda_{B,k}(u_{E_B}-v_{E_B}).
\end{equation}
Finite bounds force $h^{B,k}=0$ at zero weight. The formulation needs no
division by a variable; division appears only in the recovery proof, at
strictly positive weights. If $J_B=\varnothing$, the block has one merged
state of weight one, and its conditions reduce to $x_{E_B}-v_{E_B}\in K_B$,
which $x\in\Flow$ already enforces. Unobserved blocks therefore need no new
variables or constraints.

The theorem is a constructive specialization of the known shared-simplex
gluing principle. All blocks use the same global $y$, not independent
simplex vectors. Observed state slices are not merged, because their
domains can differ once the products are fixed.

\begin{example}[A common matrix is not sufficient for merging]\label{ex:merger-boundary}
Let $M=(-1,1,2)^\top$, $c^1=(0,1,100)^\top$, and
$c^2=(0,100,2)^\top$. Both polytopes
$Q_i=\{q\in\R:Mq\le c^i\}$ equal $[0,1]$.
For weights $\alpha_1=\alpha_2=1/2$, the weighted Minkowski sum is still
$[0,1]$, whereas the aggregated rows
$Mq\le(c^1+c^2)/2$ allow $q=2$. The two upper bounds become
$q\le101/2$ and $2q\le51$. Thus replacing state-specific systems by
the sum of their right-hand sides is generally inexact, even for a common
matrix and nonempty bounded domains; this is the standard reaggregation
boundary discussed by \citet[Section~2, equation~(7)]{KisHorvath2022}.
The two sets are identical, so the set identity
\eqref{eq:homothetic-merger} still holds; what fails is adding their
different, redundant right-hand-side descriptions. In
\cref{thm:block-state-reduction}, every state uses the single fixed
representation \eqref{eq:block-domain}, with bounds scaled by the state
weight as in \eqref{eq:linear-scaled-block}. Adding these scaled bounds gives
the same representation scaled by the total weight, so the set identity
justifies the merged formulation.
\end{example}

Finally, these results concern the hull of \eqref{eq:graph-set}. Intersecting
$\Hull$ with additional linear side constraints yields a valid relaxation
of the graph with those constraints, but need not give its exact convex
hull. Such constraints can couple blocks or alter conditional state domains;
exactness then requires a separate argument.
